\documentclass[10pt,a4paper]{amsart}
\usepackage[margin=19mm]{geometry}
\usepackage{amsmath,amssymb,mathtools}
\usepackage[scr=rsfs,cal=euler]{mathalfa}
\usepackage{xfrac}
\usepackage[unicode,hidelinks,bookmarks=false]{hyperref}
\newcommand{\ie}{i.e.\ }
\newcommand{\cf}{cf.\ }
\newcommand{\resp}{respectively\ }
\newcommand{\Subsection}{\subsection}
\providecommand{\nobreakdash}{\nobreak\hbox{-}}
\setlength{\emergencystretch}{2em}
\multlinegap0pt
\usepackage{bm}
\usepackage{bbm}
\usepackage{dsfont}
\usepackage{xspace}
\usepackage{cancel}
\usepackage{mathtools}
\usepackage[shortlabels]{enumitem}
\usepackage{amsthm}
\makeatletter
\@ifundefined{thm}{\newtheorem{thm}{Thm}}{}
\@ifundefined{lem}{\newtheorem{lem}[thm]{Lemma}}{}
\makeatother
\newcommand{\Kbar}{\overline{K}}
\newcommand{\Mod}[1]{\ \left(\operatorname{mod}\ {#1}\right)}
\newcommand{\bbP}{\mathbb{P}}
\newcommand{\calO}{\mathcal{O}}
\newcommand{\calT}{\mathcal{T}}
\newcommand{\frakP}{\mathfrak{P}}
\newcommand{\frakp}{\mathfrak{p}}
\newcommand{\frakq}{\mathfrak{q}}
\newcommand{\ord}{\operatorname{ord}}
\title[Cubic points on dynamical modular curves]{Cubic points on dynamical modular curves}
\author{John R. Doyle, Alexander Galarraga}
\date{Source: arXiv:2511.11491v1 (CC BY 4.0)}
\begin{document}
\maketitle

\noindent\emph{Source and licence.} Cubic points on dynamical modular curves. Version arXiv:2511.11491v1 (CC BY 4.0), distributed under the Creative Commons Attribution 4.0 International licence, as recorded in the arXiv licence metadata for this exact version and shown on the abstract page https://arxiv.org/abs/2511.11491v1. Full rights notice: licenses/arxiv-2511.11491-CC-BY-4.0.txt.\par\smallskip

\noindent\textit{Editorial scope.} Counted unit: the Lem whose source label is \texttt{lem:HIT} in the article (source0.tex lines 893--900, source label \texttt{lem:HIT}), delivered together with the article's own complete original proof of it (source0.tex lines 902--933, one \texttt{proof} environment of 32 lines). No mathematics is written, repaired or re-derived: statement and proof are the article's own text line for line. Required context retained uncounted: none. Every definition, display and label of those lines is retained unchanged. Omitted from this package and not redistributed: 1--892, 901--901, 934--3233. Nothing inside a retained excerpt is omitted or altered: every retained body is delivered exactly as the article's lines are written, so no omission receipt of any kind accompanies this package. Citations: the retained excerpts carry no citation command at all, so no bibliography entry of the article is transcribed and no reference is redistributed. Reference adaptation: none was needed and none was made; every reference inside the delivered unit resolves inside it, so no unresolved reference or citation is produced. Printed slips kept literal: no printed word, symbol, hypothesis or number of the counted statement or of its proof was altered. Layout and class: the article's own class is \texttt{amsart} with packages including amssymb, amsmath, caption, color, comment, dashbox, enumitem, fullpage, hyperref, longtable, mathrsfs, multirow, stmaryrd, tabularx; the wrapper is the lane's pinned \texttt{amsart} editorial wrapper at 10pt on A4, which restates the theorem-like environments the retained text needs and adds the article's own macro definitions that the retained text needs (\texttt{\textbackslash Kbar}, \texttt{\textbackslash Mod}, \texttt{\textbackslash bbP}, \texttt{\textbackslash calO}, \texttt{\textbackslash calT}, \texttt{\textbackslash frakP}, \texttt{\textbackslash frakp}, \texttt{\textbackslash frakq}, \texttt{\textbackslash ord}), with extra wrapper packages (bm, bbm, dsfont, xspace, cancel, mathtools, enumitem). The delivered numbering is the unit's own. Assets: no retained span contains a figure environment, a graphics-inclusion command, a PDF-page inclusion, a table or a TikZ picture; no image file is copied into this package, the package declares no asset dependency and no image file is redistributed with it. Licence: version arXiv:2511.11491v1 of this article is distributed under the Creative Commons Attribution 4.0 International licence, as recorded in the arXiv licence metadata for this exact version and shown on the abstract page https://arxiv.org/abs/2511.11491v1. A full rights notice is delivered at licenses/arxiv-2511.11491-CC-BY-4.0.txt. Characters: every retained excerpt is pure ASCII, so the delivered engine, XeLaTeX with its default Latin Modern text font, reports no missing character.
\begin{lem}\label{lem:HIT}
    Let $X/K$ be a curve, let $\varphi$ and $\psi$ be two dominant morphisms to $\bbP^1$, and let $d = \deg\varphi$. Then there is a finite set $S$ of maximal ideals in $\calO_K$ for which the following is true: For all maximal ideals $\frakp \notin S$, and for any finite collection of curves $X_1,\ldots,X_n$ with dominant morphisms $\pi_i : X_i \to X$, there are infinitely many points $P\in X(\Kbar)$ such that
        \begin{enumerate}[label={\textup{(\roman*)}}]
        \item $[K(P) : K] = d$,
        \item the fiber $\pi_i^{-1}(P)$ is irreducible over $K$ for all $i = 1,\ldots,n$, and
        \item for all maximal ideals $\frakq\mid\frakp$ in $\calO_{K(P)}$, we have $\ord_\frakq(\psi(P)) \ge 0$.
        \end{enumerate}
\end{lem}

\begin{proof}
Let $\pi : X \to Y$ be a morphism of maximal degree such that both $\varphi$ and $\psi$ factor through $\pi$, and let $\varphi',\psi' : Y \to \bbP^1$ be such that
    \[
        \varphi = \varphi' \circ \pi \quad\text{and}\quad \psi = \psi' \circ\pi.
    \]
Let $e, d' \ge 1$ be the degrees of $\pi$ and $\varphi'$, respectively; in particular, $d = d'e$.

Let $x$ and $y$ be the elements of the function field of $Y$ corresponding to the maps $\varphi'$ and $\psi'$, respectively. Then, since $\varphi'$ and $\psi'$ do not factor through a common map of degree at least $2$, the curve $Y$ has a plane model defined by the vanishing of a polynomial of the form
    \[
        F(x,y) = a_{d'}(x)y^{d'} + a_{d'-1}(x)y^{d'-1} + \cdots + a_1(x)y + a_0(x) \in K[x][y].
    \]

Choose $x_0 \in \calO_K$ such that $a_{d'}(x_0) \ne 0$, and let $S$ be the set of maximal ideals $\frakP \subset \calO_K$ such that either $\ord_\frakP(a_{d'}(x_0)) \ne 0$ or one of the coefficients of at least one $a_i(x) \in K[x]$ has negative $\frakP$-adic valuation. Note that $S$ is a finite set. Now let $\frakp$ be any maximal ideal not in $S$, and let
    \[
    [x_0]_\frakp := \{t \in K : t\equiv x_0 \Mod \frakp\}
    \]
be the residue class of $x_0$ modulo $\frakp$.

By Hilbert Irreducibility, if $\calT \subset K$ is a thin set (in the sense of Serre), then $[x_0]_\frakp \smallsetminus \calT$ is infinite.
Thus, if we let $\pi_i : X_i \to X$ with $1\le i \le n$ be a collection of maps as in the statement of the lemma, then the set
    \[
    \Sigma := \{t \in [x_0]_\frakp : (\varphi\circ\pi_i)^{-1}(t) \text{ is irreducible for all } i = 1,\ldots,n\}
    \]
is infinite. We claim that every point in the infinite set $\varphi^{-1}(\Sigma) \subset X(\Kbar)$ satisfies the three conditions of the lemma.

Thus, let $P \in \varphi^{-1}(\Sigma)$. For every $i = 1,\ldots,n$ and every preimage $Q_i\in\pi_i^{-1}(P)$, we have
    \[
        [K(P) : K] = \deg \varphi = d \quad\text{and}\quad [K(Q_i) : K(P)] = \deg\pi_i,
    \]
so that conditions (i) and (ii) are satisfied.
Finally, write $t = \varphi(P)$ and $u = \psi(P)$, so that $F(t,u) = 0$. Note that $t \in [x_0]_\frakp$ by assumption, so each $a_i(t)$ is integral at $\frakp$ and $a_{d'}(t)$ is a $\frakp$-adic unit. In particular, this implies that $u$ is integral over each maximal ideal $\frakq \subset \calO_{K(P)} = \calO_{K(u)}$ lying over $\frakp$, so condition (iii) is satisfied as well.
\end{proof}

\end{document}
